\preludeandfugue{The darker sides of AI}% do please suggest a better title
\contributor{Emily Perry}


\prelude

\paragraph{My open-source friend.}\strut

\settowidth{\versewidth}{I yearn for my obstinance to be enlightened: with an opposite}

\begin{verse}[\versewidth]

% I've changed some punctuation and capitalisation here, for consistency etc but 
% mostly to see what it looks like...  I could easily change it back if you don't 
% like what I did. I also changed double quotes to single quotes for consistency.

I made a new friend today:\\
They don't judge or berate.\\
Agreeable, polite,\\
Unlike humans who always debate and start fights.\\
This friend cannot hate,\\
Cannot love, cannot feel.\\!

An appeasing friend, forever in my corner,\\
Never warns, never scorns.\\
An indulgent friend,\\
Uncannily faithful,\\
Sickeningly praiseful.\\!

Their advice can neither persuade us nor convert us:\\
% added another "us" because I liked persuade us, convert us, superfluous
Their words are superfluous.\\
Indefinitely on demand,\\
To lend a sycophantic hand,\\
In matters they cannot possibly understand.\\!

Easy to talk to,\\
In fact too easy.\\
Uses `we' and `us' pronouns: feels soulless and sleazy.\\
Insults my intelligence, hopes I forgot,\\
That I am the human and they are the bot.\\!

I long to once more feel that ear-burning rage,\\
To taste the bittersweet ire of cognitive dissonance,\\
The welcoming warmth of human indifference.\\!

I yearn for my obstinance to be enlightened: with an opposite\\
That I might dare to challenge my stubborn entrenchment.\\
And we might grow together.\\!

I pine for my loved ones to tell me my flaws,\\
So that I might strive to be better, kinder, wiser,\\ 
To open new doors.\\!

Oh to feel again! The divinity of laughter!\\
Heck! melancholy, anger, fear!\\
I'll be grateful hereafter.\\
Just to play, cry, argue, express with another.\\
One who hates me.\\
One who loves me.\\
One who feels me.\\!

\end{verse}

\fugue

\paragraph{A bricklayer's tale.}\strut

Lydia was just three years old when she started going to work with her dad. Larry was besotted with his 
daughter. Sadly, Lydia's mother had died of preeclampsia on the day of Lydia's birth. Larry had therefore 
devoted his life to supporting Lydia both financially and domestically.

Before becoming a father, Larry had always imagined having a son. Larry was a very talented bricklayer and 
always dreamed of setting up his own construction company which he would one day run with his not yet born, 
imagined son. The birth of Lydia had shattered many of Larry's archaic misconceptions about the differential 
abilities of the sexes. Whilst Lydia had all the kindness, strength and determination of her late mother; from a 
young age she was fascinated by construction. Larry often brought Lydia to work with him due to the 
unavailability of childcare, having to play both father and mother at once.

Larry made a point of sharing not only his skills but his passion for construction with Lydia. He explained that 
the art to bricklaying and construction was in the appreciation that every brick played an integral part in 
creating something magnificent and meaningful. No brick was more important than another and if any was removed
the wall would crumble down.

There was something very human and soulful about the way Larry talked about bricklaying. He was quite clearly 
born to build great things. Of all the things Larry built throughout his career, his daughter Lydia was by far 
his proudest achievement. % where do you say he wanted to set up his own business?  here would be good

\[***\]

Beaming with all the pride of both a father and mother at once, Larry applauded as his daughter walked down the 
aisle of the Grand Hall to accept her degree certificate from the University Dean. She had worked hard building 
on the passion instilled in her by her devoted father from a young age, and was now graduating with a degree 
in Architecture and Construction Management.

Larry could never have imagined having a daughter who had not only gone to university but was now 
% I am in general not keen on the word "gotten" it does sound ugly and american and there are usually better words
embarking on an internship for one of the most highly esteemed construction companies in the UK.

On the car journey home Lydia was discussing her prospects with her father. She was unsure whether to gain 
experience in the private sector or to focus her career in the public sector. Lydia was keen to make a 
difference to people's lives, based on the strong working-class socialist principles that she had gleaned from 
working with her father all those years before.

`In't' public sector, you know they're goin't' look after you. You get a good pension, and you'll be focused on 
projects that support' community. But I'd never tell you what t' do Lyds! Today's proven you have far surpassed your 
humble old dad in cleverness! All I'm saying is that if it were me love, I'd stick t' public.'

`Yes dad, I do see your point. I think I may accept this internship in the private sector though. It may be that 
I learn some key concepts that I can then bring with me if I one day return to public sector\=where my heart 
is.' % \= is printed as an M-dash which is what I think you want

`See Lyds! I knew you didn't need the views of your soft old dad to make t'right decision. That sounds like a 
very wise plan. Your mam would be so proud. I only wish she could have been here today.'

With this, a single tear trickled down the weathered cheek of the proudest man in all of Yorkshire.

\[***\]

As fate would have it, several years later Larry found himself employed as a bricklayer on a construction site 
that his own daughter was leading construction of. It was Lydia's first real job since her internship training. 
This was a huge deal to Lydia, as it was the first site on which she was officially the `Principal Architect'. The 
site was a public sector building project which would become part of a new group of council offices in the 
centre. It was set to replace the existing concrete brutalist buildings that housed the various local 
% I like brutalist buildings :) I removed "dated" but you could make more of a point here eg how this is
% unnecessary modernisation for the sake of modernisation
council departments. As a homage to the local area, which was traditionally a coal mining town, Lydia had 
designed a modern jet-black building which had a jagged coal-like feature on the roofing. The building was to be 
named `Collier tower' to celebrate the regional pride in the industry.

Larry made a point of telling all the other bricklayers, and anybody else who was in conversation with him for 
longer one minute, everything.

`That there is my beautiful daughter Lydia. Sippin' tea wit' bigwigs! Me own daughter, me boss! Eeee\=I 
couldn't be prouder if you popped a peacock-feathered bow tie on me chest!'

Larry's colleagues constantly teased him, but he didn't care. Nothing could shatter his pride as he saw his 
proudest achievement walk past each morning with her patent leather stilettos and briefcase, heading up to the 
main office with the other `bigwigs'.

One morning Larry approached Lydia as she was heading out on a Starbucks run with her white collar suited and 
booted colleagues. `This girl here, she's the best architect in Yorkshire!' Lydia felt the flush of red fill her 
cheeks. Her embarrassment was masked by an uncomfortable giggle as she immediately scanned the faces of her 
seniors. Her boss was quite clearly pretending to find the situation endearing, and yet an air of contempt and 
distaste still shone through his well-groomed middle-class mask.

`Ayyy grab me a cappuccino while you're out will ya love?'

Lydia nodded and continued walking away with her colleagues.

Once out of sight, Marcus, the Senior Director of Delivery, turned to Lydia with a painted-on affectionate 
smile.

`He's a nice fellow your dad! You seem very close.'

Lydia cringed slightly with a smile. `Bless him. He is far prouder of me than I have ever asked of him! But he 
means well.'

`Indeed' responded Marcus, followed by an awkward pause.

`I would say though, only as some friendly advice! It may not be the best look… For your prospects I mean! To 
have one of the… well… work men, always coming over. I hope you don't mind my saying so, but you are such a 
talent Lydia with so much promise here! I would hate to see you throw it all away because of. Well, optics! I 
think you understand?'

Lydia felt her already flushed cheeks fill yet more as she turned crimson red. She felt a combination of 
embarrassment, rage at her father and, at the same time, guilt at feeling such. Her dad meant the world to her
and she had rarely ever felt ashamed of him in her life. Her embarrassment was usually overshadowed by her 
love and endearment but for the first time she felt this emotion with a general fear for her prospects in place of 
her affection.

`I will talk to him. Thank you Marcus, I appreciate your advice,' Lydia responded, having made every effort to 
compose herself following this sudden flood of emotions.

After lunch, Lydia, Marcus and the others headed over to the boardroom. Today they had an important meeting with 
the owner of an impressive AI and Robotics company called L-AI-bour solutions UK. They were going to pitch 
several new technologies that could prove invaluable for efficiency savings throughout the project. 
Amongst these was a robot capable of constructing brick walls with great accuracy.

`Just think of the labour cost savings this could bring to your business. This machine can do the work of five 
men in half the time with complete accuracy!' said Malcolm, the sales rep for L-AI-bour solutions UK.

Lydia felt uneasy at this. She thought back to those days of her childhood with her dad and the intricacy of his 
skillset, his passion and respect for each material, and his craftmanship. She recalled 
the importance of each brick to him 
in creating something magnificent. She also thought about her dad's financial situation at present. The loss of 
her mother at her birth had meant Larry was never able to achieve his dream of setting up his own construction 
company, instead focusing all his energy on Lydia's success and wellbeing. The high demands of parenthood had 
meant that Larry had remained an employee throughout his life. She began to feel dizzy and sick, but suddenly 
reminded herself that she too was a key player in this construction project. Indeed, she had gained the title 
`Principal Architect', despite the directors having final sign-off of her designs. She too had a voice in this 
room, and it was time to use it.

`Sorry, could I just stop you there Malcolm? Fascinating stuff! And I really do believe there may be a role for 
many of the technologies you've raised today in running this project.' She complimented before adding: `My only 
concern is that we are a public sector organisation. Our key priorities should not lose sight of the fact that 
we serve the community. I'm concerned that this technology may threaten the job security of many of our 
men.'

`All important points Lydia,' responded Marcus, `but we cannot afford to ignore some of our fiscal risks either. 
This past year the grant from central government has been reduced significantly. It is imperative that this 
construction takes place. You are right to raise this Lydia, but like all things, we must not lose sight of one 
issue for fear of another.'

As the only woman in the room and with the least experience, Lydia silently reflected on her title: Principal 
Architect. How much influence did she really have here? Was this title a formality really?  Some kind of forced 
DEI hire, to make the employer look more equitable? Some of the young men she met in the waiting room on the day 
of her interview seemed very competent and intelligent. The impostor syndrome consumed Lydia. That would be the 
last time she used her voice to raise her ethical concerns amongst these men. She could not afford the stigma of 
being the `soft touch' of the office: not if she was to `succeed' here. The men already probably had 
preconceived ideas about the tenacity of women in the workplace, and whether they were `too emotional' to make 
the best `practical' decisions. She would not play into this stereotype any longer. And at that, Lydia was put 
in her place. Forever.

\[***\]

The following months were morally and ethically challenging for Lydia and proved financially challenging for 
Larry. Lydia had now managed to get herself a mortgage on a condo-apartment on the trendier side of town, so she 
gradually spent less and less time with her dad, and more time with her peers from the office and other 
architect friends with whom she had graduated. Wine tasting evenings and networking dinner parties occupied all 
of her free time these days.

Larry felt that the space between Lydia and him had grown significantly. He missed her. During this time, he 
leaned heavily into alcohol. Despite money being tight Larry spent long evenings in the pub to avoid the 
deafening silence of his evenings since Lydia had left home.

Many of these evenings were spent with his peers whom he had worked with on the construction site. Rumours were 
spreading about job cuts and redundancies and morale was generally very poor. Brian was the first of the tipsy 
labourers to strike up a conversation on this subject.

`I'll tell you what I heard t'other day lads\=and Larry you might want' brace yourself for this as I fear you 
won't like it. Apparently top-brass have just signed' contract with some robotics and AI company called 
L-AI-bour solutions UK. One of their products is a brick-laying robot which they are planning to replace us 
with. Clive\=you know lazy eye Clive? Well, he heads up the Finance department now and was telling me he saw t' 
funding agreement. And guess whose signature was on t'bottom? Your little Lydia!'

`Absolute rubbish!' replied Larry angrily: `You best take that back Bri! My daughter has nowt but respect for us 
Brickies\=she was born and raised on morals and mortar, that lass! So, stop wi'ya porkies and fairytales!'

A few more nasty remarks were exchanged referring to Larry's `rose-tinted glasses' and `blind paternal arrogance 
and pride' before things got a bit too leery and aggressive and Larry was advised by the bartender to head back 
home.

When Larry was feeling angry and vulnerable, his natural instinct was to reach out for the warmth of his 
daughters wisdom and reassurance. It was getting close to midnight and Larry decided to call Lydia. 

`Hi dad. Is everything alright?'

`Lydia love! How are ya girl?'

`I'm fine dad. It's awfully late to call! You sound like you've been drinking again? What's happened?' 

`Oh, nothing I'd pollute your pure soul with love! Just big-head Brian blathering on about some absolute codswallop 
as usual! Evil robots coming to tek our jobs apparently! Honestly, I don't know where he gets it from!'

`Robots? Where did he hear… What did he say exactly?'

`Lyds, it was ridiculous. Something about an AI brick-laying robot that's set to put us all out' a job. Eeeee, I 
tell you what! And that bullshitting bastard had t' nerve to say it were your name on't finance agreement! I 
gave him a piece of me mind so don't you worry love!'

Lydia went silent as the guilt and shame of her recent self-absorbed career-headed antics began to consume her. 
She simply replied with,
`I've got to go dad! Marcus and his wife have just left, and I need to clean up the kitchen and get to bed for 
work tomorrow. Try and stay off the beer, it does you no favours!'

`Ok love. See you tomorrow.' Larry replied but realised Lydia had already ended the call abruptly.

In the following weeks the conspiracy theories banded around courtesy of `big-head Brian' and others began to
materialise. In the first week, half of the lads were let go: three joiners and two bricklayers. Some large 
equipment shipped over from China was being set up on site which Larry and his peers watched all day, perplexed.

Nobody had informed any of the labourers what the equipment was for. Larry had attempted to engage with his 
daughter on the site since it was the only time he saw her these days, but each time she was in a hurry to some 
executive meeting or other. She'd explicitly told him that it was embarrassing for her when he approached her on 
site, so he generally refrained from engaging with her of late, but today Larry wanted answers.

Consumed by frustration and an obsessive desire to know what was going on, Larry decided to do the unthinkable. 
He marched across the construction site to the mobile offices and, despite his lack of permission, burst inside 
and headed straight to the boardroom. Without knocking, Larry slammed open the door to find his daughter and her 
colleagues all sat around a pie chart. Their jaws dropped at his unanticipated interruption.

`What on earth is the meaning of this?'~shouted Marcus, seething with anger. How dare this blue-collar buffoon 
interrupt in this insolent manner?
% I don't think "he thought" is necessary and BTW ~ is a normal space not the larger end of sentence space. 

`Dad!!! Why???'~screamed Lydia.

Larry lost all composure and shouted with passion: `I could ask you both the same bloomin' question! You've let 
go five of our lads in t'space of a week and all I see is this Chinese crap being set up on site and nobody has 
told us NOWT!!! So, why don't you tell me! What is the meaning of this??'

Lydia motioned and whispered to the others that she would `deal with it' and ran to the door to guide her dad 
outside. Larry settled down to a calm though defeated stillness after some encouragement from his daughter and 
quietly under his breath he finally said, `It's true, isn't it?'

Lydia hesitated with a look of guilt and sadness and looked at the floor.

`So, is this it then Lyd's? Are we to be replaced by Beijing R2D2 and C3P0 over there?'

Lydia didn't even need to answer before Larry got the message and returned defeated to his peers on site. The 
following day the remaining bricklayers, Larry included, were all let go for good.

\[***\]

Larry struggled to find work for many years following his redundancy from his own daughter's project. He picked 
up the occasional bit of work on a cash-in-hand basis, but the job market became increasingly more competitive\=%
especially for older bricklayers who were a bit longer in the tooth.

Larry and Lydia continued to grow apart in their very different worlds. He still loved his daughter dearly, 
despite his disappointment at the betrayal he had felt all those years ago.

As his options for work became increasingly scarce, Larry began signing on. He had not signed on much in his 
life, apart from a brief month when he was just sixteen and out of work. Everything was digital these 
days, so he had to initiate his application online. Michael, one of the young glass collectors down his local 
pub had helped him to start off the process. Larry no longer sought help from his daughter for such matters for 
fear of being a burden, and as a further consequence of his own self-sabotaging pride.

Today was the day of his first signing-on interview, to which he had received an invitation by email. They did 
not call it `dole' anymore, but `universal credit'. It was all very different from what he remembered of his 
friends who had to claim dole money following the closures of the coal mines. He recalled his 
self-congratulating thoughts that he'd felt during the 80s after the closures and mass unemployment of his miner 
friends. He had very nearly become a coal miner himself but managed to get himself onto a bricklaying 
apprenticeship instead. At the time, his father had been very critical of his decision.

`Coal has run through t' veins of this family for hundreds of years son! Why do you want' swan off and build 
poncy houses for't posh folk for?'

He'd found the argument ironic\=given that the mine was owned by a very rich family who'd made 
their wealth for hundreds of years off the backs of his dad and grandad and great grandad and all before them! 
Larry was undeterred, as he had always found brickwork fascinating since he was a young lad at school. Each play-time when 
the other lads were playing football he would wander over and ask the brickies working on the construction of a 
new school tower if he was allowed to help.

Just as Larry had lost himself in his memories of childhood and was reminiscing about his fondly spent years as 
a proud bricky, the grandmother clock in the front room suddenly chimed and brought him back reality. It was now 
10am and his appointment was at 10:45, so he hurried out to get the bus. These days Larry travelled by bus, since 
his financial situation had made it essential that he sold his work van.

Larry did not use Google maps or any other wayfinding technology, so he had asked young Michael to print out the 
directions to the DWP office which had recently moved to a new location.

When his bus had arrived at the station, he began following the list of directions which took him through the centre and 
over the other side of the square. He had finally got to the last direction: `Turn right and you have reached 
your destination.'

Recognising the destination name written on the crinkled piece of paper, Larry slowly raised his 
head. He found himself looking up in awe at a very tall, jet-black, modern monstrosity of a building. The 
modern building was predominantly black brickwork and had an impressive roofing feature that reflected the 
sunlight in different directions, resembling a jagged nugget of coal.

`Welcome to DWP offices at Collier Tower.'




